\section{Lecture 5 (September 15th)}
\begin{defi}
(Degenerate pertubation theory) In the prior case, we've seen energy levels being strictly distinct. What if there are multiple states per energy level? 
\[P_{0}=\sum _{D}|l_{0}\rangle \langle l_{0}|\]
where $D$ is a degenerate subspace. The identity can thus be written as
\[{\bf 1}=P_0+P_1\]
Obviously,
\[\hat{H}_{0}P_0=\hat{H}_{0}\sum _{l}|l_{0}\rangle \langle l_{0}|=E_{D}^{(0)}P_0\]
and
\[P_0\hat{H}_{0}=\sum _{l}|l_0\rangle \langle l_0|\hat{H}_{0}=E_{D}^{(0)}P_0\]
and the two operators commute. We ultimately want to solve
\[\hat{H}|l\rangle =(\hat{H}_{0}+\lambda \hat{V})|l\rangle =E|l\rangle \]
we state that it can be expressed as
\[|l\rangle =P_0|l\rangle +P_{1}|l\rangle \]
Applying the Hamiltonian, we expect
\[(\hat{H}_{0}+\lambda \hat{V})(P_0|l\rangle+P_1|l\rangle )=E(P_0|l\rangle )+E(P_1|l\rangle ) \]
Applying once more the projection operator $P_0$,
\[(E_{D}^{(0)}+P_{0}\lambda \hat{V}P_0)P_0|l\rangle+P_0\lambda \hat{V}P_1|l\rangle=EP_0|l\rangle \]
we can of course subtract the most-left term and arrive at
\[(P_0\lambda \hat{V}P_0)P_0|l\rangle +P_0\lambda \hat{V}P_1|l\rangle =(E-E^{(0)}_{D})P_{D}|l\rangle \]
where the energy difference on the right is called the energy spitting $\delta E_0$. Ignoring the second term as we want only the 1st order, we have
\[(P_0\lambda VP_0)P_0|l\rangle =\delta E_{0}P_0|l\rangle \]
We can also apply $P_1$ from the left to prove that indeed the second term is of second order and beyond. Having
\[\begin{align*}
	P_1\lambda \hat{V}P_0|l\rangle +(\hat{H}_{0}+P_1\lambda \hat{V}P_1)P_1|l\rangle =&EP_1|l\rangle \\
	(E-\hat{H}_{0}-P_1\lambda VP_1)P_1|l\rangle =&P_1\lambda VP_0|l\rangle 
\end{align*}\]
which leads to
\[P_1|l\rangle =\dfrac{1}{E-H_0}P_1\lambda VP_0|l\rangle \]
\end{defi}
\vspace{2ex}
\begin{thm}The following
\[\boxed{(P_0\lambda VP_0)P_0|l\rangle =\delta E_0P_0|l\rangle }\]
is ultimately the eigenvalue problem we want to solve to find the perturbed eigenkets. The energy is resultantly
\[\boxed{E=E_0^{(0)}+\delta E_0+\lambda \sum _{k\notin D}\dfrac{|V_{kl}|^2}{E_{D}^{(0)}-E^{(0)}_{k}}}\]
\end{thm}
\vspace{2ex}
\begin{ex}
(Isotropic rigid rotator) Consider the Hamiltonian 
\[\hat{H}_{0}=\dfrac{\vec{L}\cdot \vec{L}}{2I}\]
We know that the eigenkets and their energy levels are given as 
\[|l,m\rangle\qquad \&\qquad \dfrac{l(l+1)\hbar ^2}{2I} \]
If $l\ne 0$, degenerate $2l+1$. The pertubation will be given by
\[\hat{V}=-\vec{\mu }\cdot \vec{B}=-\dfrac{gqB}{2m}L_{z}\]
where $\vec{B}$ is an external magnetic field and $\vec{\mu }=gq\vec{L}/2m$. The $g$ is dicated by whether its electron spin (=2) or whether its orbitals (=1). These kind of potentials are from Zeeman coupling. Now, we see that the potential is also commutative with $L_{z}$, as $[V,L_{z}]=0$. We should utilise this! For a general state vector,
\[\langle m|[V,L_{z}]|m'\rangle =\langle m|VL_{z}-L_{z}V|m'\rangle =0\]
which implies that
\[(m'-m)\hbar \langle m|V|m'\rangle =0\]
Thus we have found the matrix element for the pertubation! We thus find
\[\dfrac{l(l+1)\hbar ^2}{2I}-x\hbar m\]
for $x>0$.
\end{ex}
\vspace{2ex}
\begin{defi}
(Stark effect of the hydrogen atom) Assume no spin. For an Hydrogen atom in an electric field in the $z$-axis, 
\[H_1=(-e)V(r)=(-e)(-Ez)=eEz\]
\end{defi}
\vspace{2ex}

